% Generated draft: requires mathematical and source review.
\documentclass[11pt]{article}
\usepackage[margin=25mm]{geometry}
\usepackage{fontspec}
\setmainfont{Latin Modern Roman}
\usepackage{amsmath,amssymb,mathtools,unicode-math}
\setmathfont{Latin Modern Math}
\usepackage{enumitem,needspace,xurl,hyperref}
\setlength{\parindent}{0pt}
\setlength{\parskip}{5pt}
\setlength{\emergencystretch}{3em}
\newcommand{\N}{\mathbb{N}}
\newcommand{\Z}{\mathbb{Z}}
\newcommand{\Q}{\mathbb{Q}}
\newcommand{\R}{\mathbb{R}}
\newcommand{\C}{\mathbb{C}}
\newcommand{\F}{\mathbb{F}}
\newcommand{\A}{\mathbb{A}}
\newcommand{\PP}{\mathbb{P}}
\newcommand{\E}{\mathbb{E}}
\newcommand{\eps}{\varepsilon}
\newcommand{\dd}{\,\mathrm{d}}
\begin{document}
\section*{UnsolvedMath ID: 600012}
\textbf{Source problem code:} catalogue entry
\subsection{Definitions and mathematical statement}
For a real vector space \(V\) of dimension \(k+1\), write \(\PP(V)=\R P^k\) and \(\PP(V^*)\) for its dual projective space, whose points represent hyperplanes in \(\PP(V)\). A projective correlation is induced by an invertible linear map \(V\to V^*\) and extends to an incidence-reversing correspondence of projective subspaces. Projective self-duality requires such a geometric correspondence, not merely an abstract self-duality of a face poset.

The supplied research summary uses the condition \(g(P)=P^*\) for a correlation \(g:\R P^3\to(\R P^3)^*\) and a projective polyhedron \(P\). The precise geometric meaning of \(P^*\) requires the chosen polyhedron category and duality conventions. For convex polyhedra one normally works in an affine chart with a proper convex realization; the source does not say that all objects in its requested extension are convex, so that restriction is not silently imposed.

The source's two research tasks are to extend known planar projectively self-dual polygon/curve results to (i) projectively self-dual polyhedra in \(\R P^3\), and (ii) self-dual polygons in higher-dimensional real projective spaces. It supplies no one proposed classification, construction theorem or dimension formula as the exact desired output. In higher dimension it also does not define the polygon-to-dual-hyperplane indexing convention or \(m\)-self-duality, so these are not replaced by a guessed duality operation.

The nested research summary reports Ch\'avez-Caliz's 2021/2023 constructions and moduli-dimension calculations for some \(m\)-self-dual \(n\)-gons, denoting the spaces \(\mathcal M_{m,n,k}\), while leaving a general dimension formula and a higher-dimensional Clebsch conjecture open. No full citation, exact cases or formulation of that conjecture is supplied.

For polyhedra, the summary records tetrahedra as elementary projectively self-dual examples and claims a pyramid construction: for every odd \(n\geq5\), a convex \(n\)-self-dual \(n\)-gon \(Q\) in a projective plane, with respect to a polarity, and an apex outside that plane give a projectively self-dual pyramid. The field labels this new elementary progress, but its proof sketch ends with an ellipsis and has no independent verification supplied. This draft preserves it as an unreviewed source claim, not a newly proved construction or a completed classification.

The registry labels the item partially solved and its research classification PARTIAL-PROGRESS. Its provenance says the 2015 source was accessed on 2026-07-29, while current-status and rights evidence remain NEEDS\_REVIEW. Those are supplied release-review qualifications, not a new literature or rights determination. The requested extensions remain a research programme until the exact categories and success criterion are fixed.
\subsection{Short English statement}
Extend projective self-duality results from planar polygons and curves to polyhedra and higher-dimensional polygons. The source reports restricted polygon results and an incomplete, unreviewed pyramid claim, but does not specify a full classification target, higher-dimensional duality convention or precise moduli problem.
\subsection{Sources}
\begin{itemize}
\item Local UnsolvedMath registry, permanent ID 600012, source code catalogue entry, ``Baker's Dozen --- Self-dual curves and surfaces,'' including its nested research summary and release/status qualifications.
\item Tabachnikov, \emph{A Baker's Dozen of Problems} (2015), Problem 5, supplied publisher HTML source, access date 2026-07-29: \url{https://amj.math.stonybrook.edu/html-articles/Files-2015-2024/14-01/}.
\item Ch\'avez-Caliz (2021/2023), as named in the supplied summary, on \(m\)-self-dual polygons and \(\mathcal M_{m,n,k}\); no complete title, venue, identifier or URL is supplied, and none is invented.
\item Fuchs--Tabachnikov, as named in the summary's planar polygon context; no full citation is supplied. The claimed pyramid extension is an unreviewed editorial research-summary claim with a truncated proof sketch, not a separate identified publication.
\end{itemize}
\end{document}
